\exercisesection{全纯函数演算}

在以下习题中，所考虑的代数除特别声明外均以 C 为基域。

\begin{enumerate}
\def\labelenumi{\arabic{enumi})}
\tightlist
\item
  设 \(n \geqslant 1\) 为整数，且 \(A = \mathcal{L}(\mathbf{C}^{n})\)。设 \(\alpha \in C\)，并设 \(x \in A\) 为一个元素，它相对于 \(C^{n}\) 的典范基的矩阵是 n 阶、以 \(\alpha\) 为特征值的 Jordan 矩阵。x 的谱化为 \(\alpha\)。
\end{enumerate}

\begin{enumerate}
\def\labelenumi{\alph{enumi})}
\tightlist
\item
  对 \(f \in \mathcal{O}(\{\alpha\})\)，计算相对于 \(\mathbf{C}^n\) 的典范基、表示 \(f(x)\) 的矩阵。
\end{enumerate}

  设 \(m \geqslant 1\) 为整数。设 \(B_{m}\) 为 \(C^{m}\) 类函数在 \(\alpha\) 的一个邻域中的芽所成的幺 Banach 代数，赋以 \(C^{m}\) 一致收敛拓扑（VAR，R2，§12）。

\begin{enumerate}
\def\labelenumi{\alph{enumi})}
\setcounter{enumi}{1}
\item
  存在一个从 \(B_{m-1}\) 到 A 的连续幺代数态射 \(\varphi(z)=x\)，其中 \(z\in B_{m-1}\) 是 C 的恒等函数在 \(\alpha\) 处的 m−1 阶芽。
\item
  不存在连续幺代数态射 \(\varphi\)，从 \(B_m\) 到 A，使 \(\varphi(z) = x\)，其中 \(z \in B_m\) 是 \(\mathbf{C}\) 的恒等函数的芽。
\end{enumerate}

\begin{enumerate}
\def\labelenumi{\arabic{enumi})}
\setcounter{enumi}{1}
\tightlist
\item
  设 A 为幺 Banach 代数，G 为 A 的可逆元群。
\end{enumerate}

\begin{enumerate}
\def\labelenumi{\alph{enumi})}
\item
  为使元素 \(x \in A\) 具有形式 \(\exp(y)\)（其中 \(y \in A\)），充要条件是它含于 G 的一个连通交换子群之中。
\item
  为使元素 \(x \in \mathbf{A}\) 具有形式 \(\exp(y)\)（其中 \(y \in \mathbf{A}\)），只需 0 属于 \(\mathbf{C} = \mathrm{Sp}_{\mathrm{A}}(x)\) 的无界连通分支。
\end{enumerate}

\begin{enumerate}
\def\labelenumi{\arabic{enumi})}
\setcounter{enumi}{2}
\tightlist
\item
  设 A 为幺 Banach 代数，G 为 A 的可逆元群，\(G_{1}\) 为 G 的单位元连通分支。
\end{enumerate}

\begin{enumerate}
\def\labelenumi{\alph{enumi})}
\item
  设 \(x \in \mathbf{A}\)，\(n\) 为整数 \(\geqslant 0\)，使 \(x^n = e\)。证明 \(x \in \mathbf{G}_1\)。（\(x\) 的谱有限。使 \(\lambda \in \mathbf{C}\) 满足 \(y(\lambda) = \lambda x + e - \lambda e\) 为可逆者所成的集合，其余集有限，从而连通。而 \(y(0) = e\)，\(y(1) = x\)。
\item
  假设 A 交换。G/G₁ 中每个异于单位元的元素都有无穷阶。（若 \(x \in G\) 且 \(x^n \in G_1\)，则由 I，第 78 页第 10 节，有 \(x^n = \exp(y)\)（其中 \(y \in A\)），从而 \(x \exp(-y/n)^n = e\)。应用 a)。）
\end{enumerate}

\begin{enumerate}
\def\labelenumi{\arabic{enumi})}
\setcounter{enumi}{3}
\item
  存在一个交换 Banach 代数，非幺且非零，使得 \(X(A)\) 是紧的。
\item
  设 A 为幺 Banach 代数，\(a \in A\)。设 G 为 A 的可逆元群，\(G_1\) 为 G 中含单位元的连通分支。A 的指数谱是使 \(\lambda \in \mathbf{C}\) 不属于 \(x - \lambda e\) 的 \(G_1\) 所成的集合（参见 I，第 79 页的命题 12）。
\end{enumerate}

\begin{enumerate}
\def\labelenumi{\alph{enumi})}
\item
  a 的指数谱是闭的；它包含 a 的谱。
\item
  a 的指数谱的边界含于 a 的奇异谱（I，第 158 页的习题 16）。
\item
  a 的指数谱是 a 的谱与 \(\mathbf{C}=\mathrm{Sp}(a)\) 的某些有界连通分支的并。
\end{enumerate}

\begin{enumerate}
\def\labelenumi{\alph{enumi})}
\setcounter{enumi}{3}
\tightlist
\item
  设 \(f\in \mathbf{C}[\mathrm{X}]\) 为多项式。\(f\) 在 \(a\) 的指数谱上的像包含 \(f(a)\) 的指数谱。一般而言，此包含是严格的（利用 Fredholm 理论，参见 III，待出版）。
\end{enumerate}

\begin{enumerate}
\def\labelenumi{\arabic{enumi})}
\setcounter{enumi}{5}
\tightlist
\item
  设 A 为交换幺 Banach 代数。
\end{enumerate}

\begin{enumerate}
\def\labelenumi{\alph{enumi})}
\item
  若 j 是 A 的一个幂等元，则 \(\exp(2i\pi j)=1\) 且 \(\exp(i\pi j)=1-2j\)。
\item
  设 \(p \in A\) 使 \(\exp(p) = 1\)。则 \(\mathrm{Sp}_{\mathrm{A}}(p)\) 由形如 \(2in\pi\)（其中 \(n \in Z\)）的有限多个点组成。
\item
  对 \(\lambda \notin \mathrm{Sp}_{\mathrm{A}}(p)\)，我们有：
\end{enumerate}

\[
\int_ {0} ^ {1} \exp (t (p - \lambda 1)) d t = (1 - \exp (- \lambda)) (\lambda 1 - p) ^ {- 1}.
\]

\begin{enumerate}
\def\labelenumi{\alph{enumi})}
\setcounter{enumi}{3}
\tightlist
\item
  由此推出 \(\mathrm{Sp}_{\mathbf{A}}(p)\) 的每一点都是预解式 \(\mathrm{R}(p,\lambda)\) 的单极点，进而存在整数 \(m \geqslant 0\)、整数 \(n_{k} \in Z\) 以及两两正交的幂等元 \(j_{k}\)（对 \(1 \leqslant k \leqslant m\)），使得
\end{enumerate}

\[
p = 2 i \pi \sum_ {k} n _ {k} j _ {k}.
\]

\begin{enumerate}
\def\labelenumi{¶ \arabic{enumi})}
\setcounter{enumi}{6}
\tightlist
\item
  a) 设 B 为复 Banach 空间，j 为 B 中满足 \(\|j\| = 1\) 的一点。对每个 \(x \in B\)，极限
\end{enumerate}

\[
\lim_{\substack{\alpha \to 0\\ \alpha >0}}\frac{1}{\alpha} (\| j + \alpha x\| -1)
\]

存在。令 \(\varphi(x)\) 为这个极限。

\begin{enumerate}
\def\labelenumi{\alph{enumi})}
\setcounter{enumi}{1}
\tightlist
\item
  设
\end{enumerate}

\[
\psi (x) = \sup_{\substack{\varepsilon \in \mathbf{C}\\ |\varepsilon | = 1}}\varphi (\varepsilon x).
\]

函数 \(\psi\) 是一个被 B 的范数所控制的半范数。我们有

\[
\psi (x) = \sup _ {f} | f (x) |,
\]

其中 \(f\) 取遍 B 的对偶中满足 \(\|f\| = f(j) = 1\) 的元素所成的集合。

\begin{enumerate}
\def\labelenumi{\alph{enumi})}
\setcounter{enumi}{2}
\tightlist
\item
  现在假设 B 是一个 Banach 代数，且 \(j = 1\) 是 B 的单位元。我们有：
\end{enumerate}

\[
\varphi (x) = \lim_{\substack{\alpha \to 0\\ \alpha >0}}\frac{1}{\alpha}\log (\| \exp (\alpha x)\|) = \sup_{\alpha >0}\frac{1}{\alpha}\log (\| \exp (\alpha x)\|).
\]

（注意 \(\log(\|\exp((\alpha+\alpha')x)\|) \leqslant \log(\|\exp(\alpha x)\|) + \log(\|\exp(\alpha' x)\|).\)。）

\begin{enumerate}
\def\labelenumi{\alph{enumi})}
\setcounter{enumi}{3}
\tightlist
\item
  利用 c)，证明 \(\|\exp(\lambda x)\| \leqslant \exp(|\lambda|\psi(x))\)。在由 \(\lambda \mapsto \lambda^{-2} \exp(\lambda x)\) 所确定的圆周上积分 \(|\lambda| = \varrho\)，推出
\end{enumerate}

\[
\| x \| \leqslant \varrho^ {- 1} \exp (\varrho \psi (x)),
\]

从而令 \(\varrho = \psi(x)^{-1}\)，得 \(\|x\| \leqslant e\psi(x)\)，其中 \(e = \exp(1)\)。

\begin{enumerate}
\def\labelenumi{\alph{enumi})}
\setcounter{enumi}{4}
\tightlist
\item
  对 B 的每个非零元素 \(x\)，存在线性形式 \(f \in \mathrm{B}'\)，使 \(f(1) = \| f \| = 1\) 且 \(f(x) \neq 0\)。
\end{enumerate}

\begin{enumerate}
\def\labelenumi{\arabic{enumi})}
\setcounter{enumi}{7}
\item
  设 A 为幺 Banach 代数，\(x \in A\)，B 为由 x 生成的 A 的子代数。为使 \(\mathrm{Sp}_{\mathrm{A}}(x)\) 是一个有限集且其所有元素都是 x 的预解式的极点，充要条件是 B 是有限维的。（为证条件必要：若 \(\mathrm{R}(x, \lambda) = \sum_{i=1}^{p} \mathrm{F}_{i}(\lambda) a_{i}\)（其中 \(a_{i} \in A\)，\(F_{i}\) 为纯量有理函数），则 \(\mathrm{R}(x, \lambda)\) 在 \(\lambda = \infty\) 处的展开式表明诸 \(x^{n}\) 是诸 \(a_{i}\) 的线性组合。）
\item
  设 A 为幺 Banach 代数，\(x \in A\)，U 为 \(\mathrm{Sp}_{\mathrm{A}}(x)\) 的一个只有有限多个连通分支 \(U_{i}\)（\(i \in I\)）的开邻域，且 \(f \in \mathcal{O}(\mathrm{U})\)。为使 \(f(x) = 0\)，充要条件是下列条件被满足：
\end{enumerate}

\begin{enumerate}
\def\labelenumi{(\roman{enumi})}
\item
  对每个使 \(\mathrm{Sp}_{\mathrm{A}}(x)\cap\mathrm{U}_{i}\) 为无限集或含有一个不是预解式极点的点的 i，有 \(f|U_{i}=0\)；
\item
  预解式的每个 \(p\) 阶极点都是 \(\geqslant p\) 的阶 \(f\) 的零点。（关于 (ii) 的必要性，按习题 8 的方式论证。）
\end{enumerate}

\begin{enumerate}
\def\labelenumi{\arabic{enumi})}
\setcounter{enumi}{9}
\tightlist
\item
  设 \(\Delta\) 为 C 中的圆盘 \(|z| \leqslant 1\)。设 A 为由在 \(\Delta\) 上连续、在 \(\Delta\) 的内部全纯、且存在可和的复数序列 \((c_{n})\) 满足
\end{enumerate}

\[
f (z) = \sum_ {n \geqslant 0} c _ {n} z ^ {n}
\]

（对一切 \(z \in \Delta\)）的函数 f 所成的空间。这时记 \(\|f\| = \sum_{n}|c_{n}|\)。

\begin{enumerate}
\def\labelenumi{\alph{enumi})}
\item
  赋以乘法 \((f,g)\mapsto fg\)，空间 A 是一个幺 Banach 代数，且 z 是 A 的一个拓扑生成元。
\item
  空间 \(\mathsf{X}(\mathsf{A})\) 典范地等同于 \(\Delta\)；若 \(f \in A\) 且 g 在 \(f(\Delta)\) 的一个开邻域中全纯，则 \(g \circ f \in A\)。
\end{enumerate}

\begin{enumerate}
\def\labelenumi{\arabic{enumi})}
\setcounter{enumi}{10}
\item
  设 A 为交换幺 Banach 代数，\(x \in A\)。设 S 为 X(A) 的一个对 Jacobson 拓扑闭的子集，且 f 为在 \((\mathcal{G}x)(\mathrm{S})\) 的一个邻域中全纯的复值函数。存在 \(y \in A\) 使在 S 上 \(Gy = f \circ Gx\)。（设 I 为诸 \(\chi \in S\) 的核的交。在 A/I 中使用泛函演算。）
\item
  设 U 为 C 中的非空开集。证明 \(\mathcal{O}(\mathrm{U})\) 上不存在使 \(\mathcal{O}(\mathrm{U})\) 成为 Banach 代数的范数。（利用 I，第 29 页的定理 1，证明 \(\mathcal{O}(\mathrm{U})\) 中的函数都将有界。）
\item
  设 \(n\in \mathbf{N}\)
\end{enumerate}

\begin{enumerate}
\def\labelenumi{\alph{enumi})}
\tightlist
\item
  \(C^{n}\) 的子集 V 是多项式凸的，充要条件是：对每个 \(x \in C^{n} - V\)，存在整函数 \(f: C^{n} \to C\) 使得
\end{enumerate}

\[
| f (x) | > \sup _ {y \in V} | f (y) |.
\]

\begin{enumerate}
\def\labelenumi{\alph{enumi})}
\setcounter{enumi}{1}
\item
  设 \(K \subset R^{n}\) 为紧集。则 K 在 \(C^{n}\) 中是多项式凸的。
\item
  设 V 为 \(C^{n}\) 的一个多项式凸紧子集，\(f \in \mathcal{O}(V)\)。则 f 的图象在 \(C^{n+1}\) 中是多项式凸的。（利用 Oka–Weil 定理。）
\end{enumerate}

\begin{enumerate}
\def\labelenumi{\arabic{enumi})}
\setcounter{enumi}{13}
\item
  设 K 为 C 的紧子集，P（相应地，Q）为在 K 的邻域中全纯的多项式（相应地，有理）函数的芽所成集合在 \(\mathcal{O}(\mathrm{K})\) 中的闭包。证明 \(\mathrm{P} = \mathrm{Q}\) 的充要条件是 \(\mathbf{C} - \mathbf{K}\) 连通。
\item
  设 U 为模为 1 的复数所成的集合。设 \(K_{1}\)、\(K_{2}\) 为 \(C^{2}\) 中由下式给出的紧子集：
\end{enumerate}

\[
\mathrm{K} _ {1} = \mathbf {U} \times \{0 \}, \quad \mathrm{K} _ {2} = \{(z, \overline {{z}}) \mid z \in \mathbf {U} \}.
\]

\begin{enumerate}
\def\labelenumi{\alph{enumi})}
\item
  证明：在 \(K_{1}\) 的邻域中的、C 上复系数多项式函数的芽所成的集合在 \(\mathcal{O}(K_{1})\) 中不稠密。
\item
  证明：在 \(K_{2}\) 的邻域中的、C 上复系数多项式函数的芽所成的集合在 \(\mathcal{O}(K_{2})\) 中稠密。
\end{enumerate}

\begin{enumerate}
\def\labelenumi{\arabic{enumi})}
\setcounter{enumi}{15}
\tightlist
\item
  设 A 为幺 Banach 代数，\(x \in A\)，k 为整数 \(\geqslant 0\)。设 U 为模为 1 的复数所成的集合。
\end{enumerate}

\begin{enumerate}
\def\labelenumi{\alph{enumi})}
\tightlist
\item
  假设 \(\mathrm{Sp}_{\mathrm{A}}(x)\subset\mathbf{U}\)，且当 \(\|\mathrm{R}(x,\lambda)\|=\mathrm{O}((1-|\lambda|)^{k})\) 趋于 1 时 \(|\lambda|\)。证明当 n 趋于 \(\|x^{n}\|=\mathrm{O}(|n|^{k})\) 时 \(\pm\infty\)。（对每个 \(\delta>1\)，证明
\end{enumerate}

\[
2 \pi \| x ^ {n} \| \leqslant \mathrm{M} \frac {\delta^ {n}}{(\delta - 1) ^ {k}} + \mathrm{M} \frac {\delta^ {- n}}{(1 - \delta^ {- 1}) ^ {k}}
\]

其中 M 与 n 和 \(\delta\) 无关。当 n 趋于 \(\delta = n/(n - k)\) 时取 \(+\infty\)；当 n 趋于 \(\delta = n/(n + k)\) 时取 \(-\infty\)。

\begin{enumerate}
\def\labelenumi{\alph{enumi})}
\setcounter{enumi}{1}
\tightlist
\item
  假设当 n 趋于 \(\|x^{n}\|=\mathrm{O}(|n|^{k})\) 时 \(\pm\infty\)。证明 \(\operatorname{Sp}_{\mathbf{A}}(x)\subset\mathbf{U}\)，且当 \(\|\operatorname{R}(\lambda,x)\|=\operatorname{O}((1-|\lambda|)^{k+1})\) 趋于 1 时 \(|\lambda|\)。（我们有 \(\varrho(x)=\varrho(x^{-1})=1\)，从而 \(\operatorname{Sp}_{\mathbf{A}}(x)\subset\mathbf{U}\)。对 \(|\lambda|<1\)，
\end{enumerate}

\[
\mathrm{R} (\lambda , x) = - x ^ {- 1} \left(1 + \lambda x ^ {- 1} + \lambda^ {2} x ^ {- 2} + \lambda^ {3} x ^ {- 3} + \dots\right).
\]

由此导出 \(\|\mathrm{R}(\lambda,x)\|\) 的一个上界估计，并对 \(|\lambda|>1\) 的情形类似地论证。

\begin{enumerate}
\def\labelenumi{\arabic{enumi})}
\setcounter{enumi}{16}
\tightlist
\item
  设 \(X_{1}\)（相应地，\(X_{2}\)）为 \((z_{1}, z_{2}) \in \mathbf{C}^{2}\) 中满足下式的元素所成的集合：
\end{enumerate}

\[
\frac {1}{2} \leqslant | z _ {1} | ^ {2} + | z _ {2} | ^ {2} \leqslant 1, \quad \text { resp. } \quad | z _ {1} | ^ {2} + | z _ {2} | ^ {2} \leqslant 1.
\]

  设 \(X = X_{1} \cup \{(0,0)\} \subset \mathbf{C}^{2}\)，\(A = \mathcal{C}(X)\)，并设 \(a_{1}, a_{2}\) 为 \(C^{2}\) 上坐标函数在 X 上的限制。

\begin{enumerate}
\def\labelenumi{\alph{enumi})}
\item
  联合谱 \(\mathrm{Sp}_{\mathrm{A}}^{2}(a_{1},a_{2})\) 等于 X。
\item
  构造两个不同的、从 \(\mathcal{O}(\mathrm{X})\) 到 A 的连续保单位元态射，它们把 \(z_{1}\) 映为 \(a_{1}\)，把 \(z_{2}\) 映为 \(a_{2}\)。（利用下述定理：对每个在 \(f\) 的邻域中全纯的函数 \(\mathrm{X}_{1}\)，存在在 \(X_{2}\) 的邻域中全纯且在 \(X_{1}\) 的邻域中与 f 相合的函数 g（参见 I，第 168 页，习题 14）。
\end{enumerate}

\begin{enumerate}
\def\labelenumi{¶ \arabic{enumi})}
\setcounter{enumi}{17}
\tightlist
\item
  设 A 为交换幺 Banach 代数。
\end{enumerate}

\begin{enumerate}
\def\labelenumi{\alph{enumi})}
\tightlist
\item
  对 A 的元素的任何有限族 I，设 S(I) ⊂ C\(^{I}\) 为 I 的联合谱。若 I′ 是 I 的扩张，则典范投射 pr\(_{I,I'}\)：从 C\(^{I'}\) 到 C\(^{I}\) 上的投影，满足 pr\(_{I,I'}\)(S(I′)) = S(I)，从而得到 \(\varphi_{II'}\)：从 \(\mathcal{O}(S(I))\) 到 \(\mathcal{O}(S(I'))\) 中的单射态射。设 \(\mathcal{O}(X(A))\) 为诸 \(\mathcal{O}(S(I))\) 关于诸 \(\varphi_{II'}\) 的归纳极限（我们用 A × N 的有限子集来标记 A 的元素的有限族）。典范映射
\end{enumerate}

\[
\varphi_ {\mathrm{I}} \colon \mathcal {O} (\mathrm{S} (\mathrm{I})) \longrightarrow \mathcal {O} (\mathrm{X} (\mathrm{A}))
\]

都是单射的。

我们将给 \(\mathcal{O}(\mathrm{X(A)})\) 赋以诸 \(\mathcal{O}(\mathrm{S(I)})\) 的拓扑的归纳极限拓扑。对每个 \(a \in \mathrm{A}\)，\(\mathrm{S}(a) \subset \mathbf{C}\) 在 C 中某邻域内的坐标函数定义了 \(z_a\) 的一个元素 \(\mathcal{O}(\mathrm{X(A)})\)。

\begin{enumerate}
\def\labelenumi{\alph{enumi})}
\setcounter{enumi}{1}
\tightlist
\item
  连续满射 \(\mathsf{X}(\mathsf{A}) \to \mathsf{S}(\mathsf{I})\) 定义了连续态射 \(\mathcal{C}(\mathsf{S}(\mathsf{I})) \to \mathcal{C}(\mathsf{X}(\mathsf{A}))\)。态射 \(\mathcal{O}(\mathsf{S}(\mathsf{I})) \to \mathcal{C}(\mathsf{S}(\mathsf{I})) \to \mathcal{C}(\mathsf{X}(\mathsf{A}))\) 定义了一个连续态射：
\end{enumerate}

\[
\varphi \colon \mathcal {O} (\mathrm{X} (\mathrm{A})) \longrightarrow \mathcal {C} (\mathrm{X} (\mathrm{A})).
\]

这个态射一般不是单射的。

\begin{enumerate}
\def\labelenumi{\alph{enumi})}
\setcounter{enumi}{2}
\item
  存在唯一的从 \(\psi\) 到 A 的连续保单位元态射 \(\mathcal{O}(\mathsf{X}(\mathsf{A}))\)，使对一切 \(\psi(z_a) = a\) 有 \(a \in \mathsf{A}\)。（利用 I，第 51 页的定理 1。）它与 Gelfand 变换的复合是 \(\varphi\)。
\item
  设 E 为分离局部凸空间，U 为 E 的弱开子集，\(f: U \to C\) 为一个函数。若对每个 \(u \in U\)，存在 u 的一个弱开邻域 \(V_u\)、E 中余维数有限的闭向量子空间 \(F_u\)、以及 \(g_u\) 上的全纯函数 \(p_u(V_u)\)（其中 \(p_u\) 表示 E 到 \(E/F_u\) 上的典范映射），使得 \(f|V_u = g_u \circ p_u\)，则称 f 是弱全纯的。
\end{enumerate}

  证明：若 \(f\) 弱全纯，且 K 是 U 的一个弱紧子集，则存在 K 在 U 中的一个弱开邻域 \(V_K\)、E 中余维数有限的闭向量子空间 \(F_K\)、以及 \(g_K\) 上的全纯函数 \(p_K(V_K)\)，使得 \(f|V_K = g_K \circ p_K\)，其中 \(p_K\) 是 E 到 \(E/F_K\) 上的典范映射。

\begin{enumerate}
\def\labelenumi{\alph{enumi})}
\setcounter{enumi}{4}
\tightlist
\item
  证明 \(\mathcal{O}(\mathsf{X}(\mathsf{A}))\) 是诸 \(\mathcal{O}(\mathsf{S})\) 的归纳极限，其中 S 取遍 A 的对偶中 \(\mathsf{X}(\mathsf{A})\) 的弱开邻域所成的集合，而 \(\mathcal{O}(\mathsf{S})\) 表示 S 上弱全纯函数所成的代数，赋以紧收敛拓扑。
\end{enumerate}

